% TOP_PROBLEM: {"id":30007640,"problem_number":"LOCAL-30007640","title":"Remote work and career development","category_id":21,"category":{"id":21,"name":"economics","display_name":"Economics","description":"Open economic theory statements, published research agendas, and proposed scoped specifications; question provenance and historical priority are qualified per record.","slug":"economics","order_index":21,"created_at":"2026-10-08T00:00:00Z"},"status":"provenance_pending","external_url":null,"legacy_ids":[],"published":false,"statement":"In the explicitly specified setting: How does remote or hybrid work affect learning, mentoring, promotion, and lifetime earnings across occupations and career stages? The mathematical statement above fixes the target, assumptions and scope of this catalogue formulation.","economics":{"original_id":129,"field":"Labor markets and work","kind":"empirical","formalization_basis":"proposed_specification","source_status":"not_verified","priority_status":"not_established","priority_note":"No qualifying problem statement was directly verified, so historical priority cannot be assigned.","earliest_verified_reference":null,"current_reference":{"authors":["Cevat Giray Aksoy","Nicholas Bloom","Steven J. Davis","Victoria Marino","Cem Özgüzel"],"title":"Remote Work, Employee Mix, and Performance","year":2025,"url":"https://docs.iza.org/dp17917.pdf","doi":"10.3386/w33851","locator":"Section 5, Conclusion, printed pp. 9–10","evidence_summary":"The study finds onboarding and productivity differences in one call center; it does not explicitly state the broader lifetime-career question."},"formalization_note":"This is a proposed scoped research specification. No primary passage explicitly stating this entire remaining question as open was verified. The supporting literature may answer important subquestions; this entry must not be advertised as a verified formal open problem."},"statusQualification":"Provenance pending: an explicit published statement of this exact open question was not verified. This is a catalogue-authored candidate specification, not a certified open conjecture. This is a proposed scoped research specification. No primary passage explicitly stating this entire remaining question as open was verified. The supporting literature may answer important subquestions; this entry must not be advertised as a verified formal open problem.","statusReviewedAt":"2026-10-08"}
% FIRST_POSTED: 2026-10-08
% ENTRY_KIND: statement_only
% !TeX program = lualatex
\documentclass[11pt]{article}
\usepackage{fontspec}
\usepackage[a4paper,margin=25mm]{geometry}
\usepackage{amsmath,amssymb,mathtools}
\usepackage{xurl}
\usepackage[hidelinks]{hyperref}
\newcommand{\sref}[2]{\hyperlink{source:#1:#2}{[S#2]}}
\setlength{\emergencystretch}{3em}
\begin{document}
% problemId:30007640
\section*{Remote work and career development}\label{30007640}
\subsection{Definitions and mathematical statement}
Consider workers hired into specified occupations at baseline and a feasible work-location policy \(r\in\{0,h,1\}\), denoting onsite, a declared hybrid schedule, and fully remote work. Let \(Y_{it}(r)\) be earnings, \(S_{it}(r)\) independently assessed skills, and \(P_{it}(r)\) promotion status at date \(t\). Let \(X_i\) include baseline experience, occupation, and commuting constraints. Define \(\tau_Y(r,x)=E[\sum_{t=0}^{T}\delta^t\{Y_{it}(r)-Y_{it}(0)\}\mid X_i=x]\), where \(\delta\in(0,1]\) is a specified discount factor and \(T\) the follow-up horizon. Estimate analogous skill and promotion effects and determine whether mentoring, contact with experienced coworkers, or evaluation practices mediate them. Initial onboarding and subsequent work location should be varied separately. Randomized location policies or credible quasi-experiments must address worker selection, endogenous exits, and effects on coworkers. Promotion differences alone cannot distinguish learning from managerial visibility. Evaluate benefits and losses for entrants separately from experienced workers. The cited call-center findings resolve some local productivity comparisons; extrapolating those results to long-term career development and other occupations is the proposed remaining empirical task.

\par\textbf{Formulation and scope.} This is a proposed scoped research specification. No primary passage explicitly stating this entire remaining question as open was verified. The supporting literature may answer important subquestions; this entry must not be advertised as a verified formal open problem.

\par\textbf{Open-question provenance.} An explicit published open-question statement was not verified. Sources below are supporting literature and must not be treated as a first listing.

\par\textbf{Historical priority.} No qualifying problem statement was directly verified, so historical priority cannot be assigned.

\subsection{Short English statement}
In the explicitly specified setting: How does remote or hybrid work affect learning, mentoring, promotion, and lifetime earnings across occupations and career stages? The mathematical statement above fixes the target, assumptions and scope of this catalogue formulation.

\subsection{Sources}
\begin{itemize}\raggedright
\item[\textnormal{[S1]}]\hypertarget{source:30007640:1}{} Cevat Giray Aksoy, Nicholas Bloom, Steven J. Davis, Victoria Marino, Cem Özgüzel. 2025. Remote Work, Employee Mix, and Performance. Section 5, Conclusion, printed pp. 9–10.
\url{https://docs.iza.org/dp17917.pdf}.
DOI: 10.3386/w33851.
{\small\textit{Supporting or current literature; see evidence qualification. The study finds onboarding and productivity differences in one call center; it does not explicitly state the broader lifetime-career question.}}
\end{itemize}
\end{document}
